\subsection{THE RING R(G)}

Let \(\mathrm{R}(\mathrm{G})\) be the ring of representations (continuous, on finite dimensional complex vector spaces) of \(\mathrm{G}\) (Algebra, Chap. VIII, §10, no. 6). If \(\tau\) is a representation of \(\mathrm{G}\) , denote by \([\tau]\) its class in \(\mathrm{R}(\mathrm{G})\) ; if \(\tau\) and \(\tau'\) are two representations of \(\mathrm{G}\) , we have, by definition,

\[
\begin{array}{r} [ \tau ] + [ \tau^ {\prime} ] = [ \tau \oplus \tau^ {\prime} ] \\ \relax [ \tau ] [ \tau^ {\prime} ] = [ \tau \otimes \tau^ {\prime} ]. \end{array}
\]

Since every representation of G is semi-simple, the Z-module \(\mathrm{R}(\mathrm{G})\) is free and admits as a basis the set of classes of irreducible representations of G, a set that can be identified with \(X_{++}\) by Th. 1. The map \(\tau \mapsto \mathrm{L}(\tau)_{(\mathbf{C})}\) induces a homomorphism l from the ring \(\mathrm{R}(\mathrm{G})\) to the ring \(\mathcal{R}(\mathfrak{g}_{\mathbf{C}})\) of representations of \(g_{C}\) (Chap. VIII, §7, no. 6).

Let \(\tau : \mathbf{G} \to \mathbf{GL}(\mathbf{V})\) be a representation of \(\mathbf{G}\) ; we consider the gradation \((\mathrm{V}_{\lambda}(\mathrm{T}))_{\lambda \in \mathrm{X}(\mathrm{T})}\) of the \(\mathbf{C}\) -vector space \(\mathbf{V}\) . Denote by \(\mathrm{Ch}(\mathbf{V})\) , or by \(\mathrm{Ch}(\tau)\) , the character of the graded vector space \(\mathbf{V}\) (Chap. VIII, §7, no. 7); if \((e^{\lambda})_{\lambda \in \mathrm{X}(\mathrm{T})}\) denotes the canonical basis of the algebra \(\mathbf{Z}[\mathrm{X}(\mathrm{T})] = \mathbf{Z}^{(\mathrm{X}(\mathrm{T}))}\) , we have, by definition,

\[
\operatorname{Ch} (\tau) = \sum_ {\lambda \in \mathrm{X} (\mathrm{T})} \left(\dim \mathrm{V} _ {\lambda} (\mathrm{T})\right) e ^ {\lambda}.
\]

We define in the same way (loc. cit.) a homomorphism of rings, also denoted by Ch, from R(G) to Z[X(T)]. If G is semi-simple, we have a commutative diagram

\[
\begin{array}{c c c} \mathrm{R(G)} & \stackrel {{\mathrm{Ch}}} {{\longrightarrow}} & \mathbf {Z} [ \mathrm{X(T)} ] \\ \Big \downarrow_ {l} & & \Big \downarrow_ {\tilde {\delta}} \\ \mathcal {R} (\mathfrak {g} _ {\mathbf {C}}) & \stackrel {{\mathrm{ch}}} {{\longrightarrow}} & \mathbf {Z} [ \mathrm{P} ] \end{array}\tag{1}
\]

where P denotes the group of weights of \(\mathrm{R}(\mathfrak{g}_{\mathbf{C}}, \mathfrak{t}_{\mathbf{C}})\) and \(\tilde{\delta}\) the homomorphism induced by \(\delta\) .

The Weyl group W operates by automorphisms on the group X(T), and hence on the ring Z[X(T)]. By Prop. 5 of §4, no. 3, the image of Ch is contained in the subring Z[X(T)] \(^{W}\) consisting of the elements invariant under W.

PROPOSITION 2. The homomorphism \(\mathrm{Ch}\) induces an isomorphism from \(\mathrm{R}(\mathrm{G})\) to \(\mathrm{Z}[\mathrm{X}(\mathrm{T})]^{\mathrm{W}}\) .

For \(\lambda \in \mathrm{X}_{++}\) , denote by \([\lambda]\) the class in \(\mathrm{R}(\mathrm{G})\) of the irreducible representation of highest weight \(\lambda\) . Since the family \(([\lambda])_{\lambda \in \mathrm{X}_{++}}\) is a basis of the \(\mathbf{Z}\) -module \(\mathrm{R}(\mathrm{G})\) , it suffices to prove the following assertion:

The family \((\mathrm{Ch}[\lambda])_{\lambda\in\mathrm{X}_{++}}\) is a basis of the Z-module \(\mathbf{Z}[\mathrm{X}(\mathrm{T})]^{\mathrm{W}}\) .

For every element \(u = \sum_{\lambda} a_{\lambda} e^{\lambda}\) of \(\mathbf{Z}[X(T)]\) , a term \(a_{\lambda} e^{\lambda}\) in u is called maximal if \(\lambda\) is a maximal element of the set of \(\mu \in X(T)\) such that \(a_{\mu} \neq 0\) . Th. 1 implies that \(\operatorname{Ch}[\lambda]\) has a unique maximal term, namely \(e^{\lambda}\) . Thus, the proposition follows from the following lemma.

Lemma 3. For each \(\lambda \in \mathrm{X}_{++}\) , let \(\mathrm{C}_{\lambda}\) be an element of \(\mathbf{Z}[\mathrm{X(T)}]^{\mathrm{W}}\) having unique maximal term \(e^{\lambda}\) . Then the family \((\mathrm{C}_{\lambda})_{\lambda \in \mathrm{X}_{++}}\) is a basis of \(\mathbf{Z}[\mathrm{X(T)}]^{\mathrm{W}}\) .

The proof is identical to that of Prop. 3 of Chap. VI, §3, no. 4, replacing A by Z, P by X(T) and P ∩ \(\overline{C}\) by \(X_{++}\) .

Let \(\Theta (\mathrm{G})\) (resp. \(\Theta (\mathrm{T})\) ) be the \(\mathbf{C}\) -algebra of continuous representative functions on \(\mathrm{G}\) (resp. T), and let \(\mathrm{Z}\Theta (\mathrm{G})\) (resp. \(\Theta (\mathrm{T})^{\mathrm{W}}\) ) be the subalgebra consisting of the central (resp. W-invariant) functions. The restriction map \(\Theta (\mathrm{G})\to \Theta (\mathrm{T})\) induces a homomorphism of rings \(r:\mathrm{Z}\Theta (\mathrm{G})\rightarrow \Theta (\mathrm{T})^{\mathrm{W}}\) . On the other hand, the map that associates to a representation \(\tau\) its character (that is, the function \(g\mapsto \mathrm{Tr}\tau (g)\) ) extends to a homomorphism of \(\mathbf{C}\) -algebras \(\mathrm{Tr}:\mathbf{C}\otimes_{\mathbf{Z}}\mathrm{R}(\mathrm{G})\to \mathrm{Z}\Theta (\mathrm{G})\) which, by Spectral Theories, is an isomorphism. Similarly, the canonical injection \(\mathrm{X(T)}\to \Theta (\mathrm{T})\) induces an isomorphism of \(\mathbf{C}\) -algebras \(\iota :\mathbf{C}[\mathrm{X(T)}]\to \Theta (\mathrm{T})\) , which induces an isomorphism \(\iota :\mathbf{C}[\mathrm{X(T)}]^{\mathrm{W}}\to \Theta (\mathrm{T})^{\mathrm{W}}\) . The diagram

\[
\begin{array}{c c c} \mathbf {C} \otimes_ {\mathbf {Z}} \mathrm{R(G)} & \stackrel {{1 \otimes \mathrm{Ch}}} {{\longrightarrow}} & \mathbf {C} [ \mathrm{X(T)} ] ^ {\mathrm{W}} \\ \Big \downarrow_ {\mathrm{Tr}} & & \Big \downarrow_ {\iota} \\ \mathrm{Z} \Theta (\mathrm{G}) & \stackrel {{r}} {{\longrightarrow}} & \Theta (\mathrm{T}) ^ {\mathrm{W}} \end{array}\tag{2}
\]

is commutative: indeed, for every representation \(\tau : \mathrm{G} \to \mathbf{GL}(\mathrm{V})\) and all \(t \in \mathrm{T}\) ,

\[
\operatorname{Tr} \tau (t) = \sum_ {\lambda \in \mathrm{X} (\mathrm{T})} \left(\dim \mathrm{V} _ {\lambda} (\mathrm{T})\right) \lambda (t) = \iota (\operatorname{Ch} \tau) (t),
\]

that is, \((r\circ \mathrm{Tr})[\tau ] = (\iota \circ \mathrm{Ch})[\tau ]\)

We can now deduce from the proposition the following result.

COROLLARY. The restriction map \(r: \mathrm{Z}\Theta(\mathrm{G}) \to \Theta(\mathrm{T})^{\mathrm{W}}\) is bijective.

